% Copyright 2018-2022 FIUS
%
% This file is part of theo-vorkurs-folien.
%
% theo-vorkurs-folien is free software: you can redistribute it and/or modify
% it under the terms of the GNU General Public License as published by
% the Free Software Foundation, either version 3 of the License, or
% (at your option) any later version.
%
% theo-vorkurs-folien is distributed in the hope that it will be useful,
% but WITHOUT ANY WARRANTY; without even the implied warranty of
% MERCHANTABILITY or FITNESS FOR A PARTICULAR PURPOSE.  See the
% GNU General Public License for more details.
%
% You should have received a copy of the GNU General Public License
% along with theo-vorkurs-folien.  If not, see <https://www.gnu.org/licenses/>.

\documentclass[a4paper, 11pt]{scrreprt}
\usepackage{paperandpencil}
\usepackage[ngerman]{babel}
\usepackage[top=2.5cm,bottom=2.5cm,left=2cm,right=2cm]{geometry}
\usepackage{graphicx}
\usepackage{array}
\usepackage{eso-pic}
\usepackage{qrcode}
\usepackage{forloop}
\newcolumntype{L}[1]{>{\raggedright\let\newline\\\arraybackslash\hspace{0pt}}m{#1}}
\newcolumntype{C}[1]{>{\centering\let\newline\\\arraybackslash\hspace{0pt}}m{#1}}
\newcolumntype{R}[1]{>{\raggedleft\let\newline\\\arraybackslash\hspace{0pt}}m{#1}}
\IfFileExists{../config.tex}{% ==============================================================================
% ZENTRALE KONFIGURATION - JAVA / PSE VORKURS
% ==============================================================================

\ifdefined\JVKCONFIGLOADED\endinput\fi
\def\JVKCONFIGLOADED{1}

% ------------------------------------------------------------------------------
% 1. ALLGEMEINE INFORMATIONEN ZUM VORKURS
% ------------------------------------------------------------------------------
\newcommand{\vkTitle}{PSE - Vorkurs}
\newcommand{\vkYear}{2026}
\newcommand{\vkMonth}{10}
\newcommand{\vkSemester}{Wintersemester 2026/27}
\newcommand{\vkInstitute}{FIUS - Fachgruppe Informatik Universität Stuttgart}
\newcommand{\vkAuthors}{Linus, Philipp, Tillmann, Tobias}
\newcommand{\vkJdkVersion}{21}

% ------------------------------------------------------------------------------
% 2. TERMINE & DATEN DER EINZELNEN TAGE
% ------------------------------------------------------------------------------
\newcommand{\vkDateDayZero}{05.10.2026}   % Tag 0 (Einrichtung / Montag Vormittag)
\newcommand{\vkDateDayOne}{05.10.2026}    % Tag 1 (Montag)
\newcommand{\vkDateDayTwo}{06.10.2026}    % Tag 2 (Dienstag)
\newcommand{\vkDateDayThree}{07.10.2026}  % Tag 3 (Mittwoch)
\newcommand{\vkDateDayFour}{08.10.2026}   % Tag 4 (Donnerstag)
\newcommand{\vkDateDayFive}{09.10.2026}   % Tag 5 (Freitag)

% ------------------------------------------------------------------------------
% 3. ZENTRALE WEB-LINKS (FOLIEN, FEEDBACK & MUSTERLÖSUNGEN)
% ------------------------------------------------------------------------------
% Hauptwebseite mit allen Folien und Aufgabenblättern:
\newcommand{\vkFolienUrl}{https://fius.de/angebote/pse-vorkurs-materialien/}

% Feedback-Formulare (z. B. Google Forms / EvaSys):
% Bleibt ein Link leer (""), wird der Feedback-Block auf den Folien nicht angezeigt.
\newcommand{\vkFeedbackUrlGeneral}{}                                    % Allgemeiner Feedback-Link
\newcommand{\vkFeedbackUrlDayOne}{}                                     % Tag 1 (optional)
\newcommand{\vkFeedbackUrlDayTwo}{}                                     % Tag 2 (optional)
\newcommand{\vkFeedbackUrlDayThree}{}                                   % Tag 3 (optional)
\newcommand{\vkFeedbackUrlDayFour}{}                                    % Tag 4
\newcommand{\vkFeedbackUrlDayFive}{}                                    % Tag 5

\newcommand{\vkUploadsBaseUrl}{https://fius.de/wp-content/uploads/\vkYear/\vkMonth}

% Download-Links für Musterlösungen und Starter-Pakete (ZIP-Archive):
\newcommand{\vkUrlDayOneHighperformer}{\vkUploadsBaseUrl/Day1Highperformer.zip}
\newcommand{\vkUrlSolutionDayOne}{\vkUploadsBaseUrl/Day1Musterloesung.zip}
\newcommand{\vkUrlSolutionDayTwo}{\vkUploadsBaseUrl/Day2Musterloesung.zip}
\newcommand{\vkUrlSolutionDayThree}{\vkUploadsBaseUrl/Day3Musterloesung.zip}
\newcommand{\vkUrlSolutionDayFour}{\vkUploadsBaseUrl/Day4Musterloesung.zip}

% Externe Dokumentations-Links:
\newcommand{\vkUrlJavaDocs}{https://docs.oracle.com/javase/8/docs/api/}
\newcommand{\vkUrlEduroam}{uni-stuttgart.de/eduroam}
\newcommand{\vkUrlJetbrainsToolbox}{jetbrains.com/toolbox-app/}

% ------------------------------------------------------------------------------
% 4. ABENDPROGRAMM / RAHMENPROGRAMM (SOCIAL EVENTS)
% ------------------------------------------------------------------------------
% Tag 0 (Folie: Gemeinsame Aktionen):
\newcommand{\vkEventDayZero}{\textbf{Karaoke} - Heute Abend ab 18 Uhr}

% Tag 1 (Übungsblatt Kopfzeile / Ankündigung):
\newcommand{\vkEventDayOne}{Heute Abend gibts Karaoke!}

% Tag 4 (Übungsblatt: Kneipentour o. Ä.):
\newcommand{\vkEventDayFourTitle}{Kneipentour}
\newcommand{\vkEventDayFourText}{Heute Abend (direkt 18:00 Treffpunkt) ist die Kneipentour und wir vier Orgas haben mies Bock, kommt gern vorbei ob ihr Trinker seid oder nicht.}

% Tag 5 (Übungsblatt: UNO-Party / Shotbar):
\newcommand{\vkEventDayFiveTitle}{Yippie UNO}
\newcommand{\vkEventDayFiveText}{Heute Abend ist UNO, kommt gerne zwischen 23:00 und 1:00 bei der Shotbar vorbei, und wenn ihr uns ein Bild von diesem ASCII Grill zeigt gibts sogar einen shot kostenlos (gibt auch Alk-freie Shots)!}

% ------------------------------------------------------------------------------
% 5. FEEDBACKBÖGEN / UMFRAGEN (PRINT SURVEYS)
% ------------------------------------------------------------------------------
\newcommand{\vkSurveyTitle}{Feedback - \vkTitle\ \vkYear}
\newcommand{\vkSurveyRooms}{$\Box$ 0.108 & $\Box$ 0.124 & $\Box$ GS-Pool & $\Box$ 0.363 \\ & $\Box$ 0.463 & $\Box$ 0.447 & $\Box$ 0.457}
}{}
\pagestyle{plain}
\begin{document}

\providecommand{\docNum}{0}
\newcounter{x}
\setcounter{x}{\docNum}
%\forloop{x}{0}{\value{x} < 50}{%
\pagenumbering{arabic}
\setcounter{page}{1}



\begin{figure}[h]
  \centering
  % \includegraphics{school_logo.png} Optional -uncomment to use
\end{figure}
\section*{\LARGE \vkSurveyTitle}
\vspace{1cm}

\begin{tabular}{p{5.6cm}p{1.8cm}p{1.8cm}p{2.3cm}p{1.8cm}p{1.8cm}p{0.001cm}}
  Ich war heute im Seminarraum... & \vkSurveyRooms
\end{tabular}

\question*{Welchen Studiengang und im wievielten Semester studierst du?}
\fbox{\parbox[c][1cm]{\textwidth}{\phantom{.}}}

%\section*{Your job}
%\question{FRAGE}
\subsection*{MINT Vorkurse}
\begin{tabular}{p{9.5cm}C{2.1cm}C{2.1cm}p{0.001cm}}
   & ja & nein                                       \\
  Ich war auch im Informatik-Vorkurs des MINT-Kolleg  &  $\Box$ &  $\Box$\\
\end{tabular}
\\

\subsection*{Allgemein}
\vertikalblockfour{stimme zu}{stimme eher zu}{stimme eher nicht zu}{stimme nicht zu}{
  \blocktextfour{Der Vorkurs hat heute Spaß gemacht.}
  \blocktextfour{Der Vorkurs war gut organisiert.}
}

\subsection*{Vortrag}
\vertikalblockfour{stimme zu}{stimme eher zu}{stimme eher nicht zu}{stimme nicht zu}{
  \blocktextfour{Die Vortragenden haben den Inhalt verständlich erklärt.}
  \blocktextfour{Der Vortrag wurde professionell durchgeführt.}
  \blocktextfour{Die Vortragenden wirkten sympathisch.}
  \blocktextfour{Die Kombination aus Vortrag, Quizen und Live-Coding war gut.}
}

\subsection*{Übungsaufgaben}
\vertikalblockfour{stimme zu}{stimme eher zu}{stimme eher nicht zu}{stimme nicht zu}{
  \blocktextfour{Die Übungsaufgaben hatten einen Mehrwert für den Vorkurs.}
  \blocktextfour{Bei Fragen habe ich schnell Hilfe bekommen.}
  \blocktextfour{Bei Fragen konnte mir gut weitergeholfen werden.}
  \blocktextfour{Ich konnte mich mit den anderen Teilnehmern austauschen.}
}

\subsection*{Tempo}
\vertikalblockfive{viel zu hoch}{zu hoch}{genau richtig}{zu niedrig}{viel zu niedrig}{
  \blocktextfive{Das Tempo in dem Inhalte im Vortrag behandelt wurden war...}
  \blocktextfive{Die Schwierigkeit der Aufgaben war...}
  \blocktextfive{Die Anzahl der Aufgaben...}
}


\subsection*{Weiterer Vorkurs}
\vertikalblocktwo{versuchen jeden Tag den Vorkurs zu besuchen.}{nicht mehr kommen.}{
  \blocktexttwo{Ich werde...}
}


\subsection*{Sonstiges}


\question*{Was hat dir besonders gut gefallen?}
\fbox{\parbox[c][3.5cm]{\textwidth}{\phantom{.}}}

\question*{Was hat dir nicht so gut gefallen?}
\fbox{\parbox[c][3.5cm]{\textwidth}{\phantom{.}}}

\question*{Welche Verbesserungsvorschläge hast du für uns?}
\fbox{\parbox[c][3.5cm]{\textwidth}{\phantom{.}}}

\question*{Was ich sonst noch sagen wollte...}
\fbox{\parbox[c][4cm]{\textwidth}{\phantom{.}}}

\AddToShipoutPictureBG*{
  \put(0.5\paperwidth+5mm,7.5mm+5mm){\qrcode[height=15mm]{\arabic{x}S4}}
}

\pagebreak

%}

\end{document}
